\documentclass[11pt]{article}
\ifdefined\pdftrailerid\pdftrailerid{}\fi
\ifdefined\pdfinfoomitdate\pdfinfoomitdate=1\fi
\def\DoNotLoadEpstopdf{}
\usepackage[T1]{fontenc}
\usepackage{lmodern}
\usepackage{mathrsfs}
\usepackage{amsmath,amssymb,amsthm,mathtools}
\usepackage[margin=1in]{geometry}
\usepackage{microtype}
\usepackage{enumitem}
\usepackage{xcolor}
\usepackage[colorlinks=true,linkcolor=blue!45!black,citecolor=blue!45!black,urlcolor=blue!45!black]{hyperref}

\ifdefined\pdfglyphtounicode
  \pdfgentounicode=1
  \pdfglyphtounicode{tfm:lmsy10/mapsto}{007C}
  \pdfglyphtounicode{tfm:lmsy10/E}{2130}
  \pdfglyphtounicode{tfm:lmsy10/H}{210B}
  \pdfglyphtounicode{tfm:lmsy10/O}{1D4AA}
  \pdfglyphtounicode{tfm:lmsy10/R}{211B}
  \pdfglyphtounicode{tfm:lmsy8/H}{210B}
  \pdfglyphtounicode{tfm:rsfs10/G}{1D4A2}
  \pdfglyphtounicode{tfm:msbm10/E}{1D53C}
  \pdfglyphtounicode{tfm:msbm10/F}{1D53D}
  \pdfglyphtounicode{tfm:msbm10/P}{2119}
  \pdfglyphtounicode{tfm:msbm10/R}{211D}
  \pdfglyphtounicode{tfm:lmex10/parenleftbig}{0028}
  \pdfglyphtounicode{tfm:lmex10/parenrightbig}{0029}
  \pdfglyphtounicode{tfm:lmex10/vextendsingle}{23D0}
  \pdfglyphtounicode{tfm:lmex10/vextenddouble}{2016}
  \pdfglyphtounicode{tfm:lmex10/parenleftBig}{0028}
  \pdfglyphtounicode{tfm:lmex10/parenrightBig}{0029}
  \pdfglyphtounicode{tfm:lmex10/parenleftbigg}{0028}
  \pdfglyphtounicode{tfm:lmex10/parenrightbigg}{0029}
  \pdfglyphtounicode{tfm:lmex10/parenleftBigg}{0028}
  \pdfglyphtounicode{tfm:lmex10/parenrightBigg}{0029}
  \pdfglyphtounicode{tfm:lmex10/braceleftBigg}{007B}
  \pdfglyphtounicode{tfm:lmex10/angbracketleftBigg}{27E8}
  \pdfglyphtounicode{tfm:lmex10/angbracketrightBigg}{27E9}
  \pdfglyphtounicode{tfm:lmex10/parenlefttp}{239B}
  \pdfglyphtounicode{tfm:lmex10/parenrighttp}{239E}
  \pdfglyphtounicode{tfm:lmex10/bracketlefttp}{23A1}
  \pdfglyphtounicode{tfm:lmex10/bracketrighttp}{23A4}
  \pdfglyphtounicode{tfm:lmex10/bracketleftbt}{23A3}
  \pdfglyphtounicode{tfm:lmex10/bracketrightbt}{23A6}
  \pdfglyphtounicode{tfm:lmex10/parenleftbt}{239D}
  \pdfglyphtounicode{tfm:lmex10/parenrightbt}{23A0}
  \pdfglyphtounicode{tfm:lmex10/parenleftex}{239C}
  \pdfglyphtounicode{tfm:lmex10/parenrightex}{239F}
  \pdfglyphtounicode{tfm:lmex10/circleplustext}{2A01}
  \pdfglyphtounicode{tfm:lmex10/circleplusdisplay}{2A01}
  \pdfglyphtounicode{tfm:lmex10/summationtext}{2211}
  \pdfglyphtounicode{tfm:lmex10/summationdisplay}{2211}
  \pdfglyphtounicode{tfm:lmex10/productdisplay}{220F}
  \pdfglyphtounicode{tfm:lmex10/hatwide}{02C6}
  \pdfglyphtounicode{tfm:lmex10/hatwidest}{02C6}
  \pdfglyphtounicode{tfm:lmex10/tildewide}{02DC}
  \pdfglyphtounicode{tfm:lmex10/radicalbt}{23B7}
  \pdfglyphtounicode{tfm:lmex10/radicalvertex}{23D0}
  \pdfglyphtounicode{tfm:lmex10/radicaltp}{231C}
\fi

\newtheorem{theorem}{Theorem}[section]
\newtheorem{proposition}[theorem]{Proposition}
\newtheorem{lemma}[theorem]{Lemma}
\newtheorem{corollary}[theorem]{Corollary}
\theoremstyle{definition}
\newtheorem{definition}[theorem]{Definition}
\theoremstyle{remark}
\newtheorem{remark}[theorem]{Remark}
\newcommand{\E}{\mathbb E}
\newcommand{\F}{\mathbb F}
\newcommand{\R}{\mathbb R}
\newcommand{\N}{\mathbb N}
\newcommand{\eps}{\varepsilon}
\newcommand{\Id}{\mathrm{Id}}
\newcommand{\GL}{\mathrm{GL}}
\newcommand{\rank}{\operatorname{rank}}
\newcommand{\Span}{\operatorname{span}}
\newcommand{\Odd}{\mathcal O}
\newcommand{\Rad}{\mathcal R}
\newcommand{\norm}[1]{\left\lVert #1\right\rVert}
\newcommand{\abs}[1]{\left\lvert #1\right\rvert}
\newcommand{\pair}[2]{\left\langle #1,#2\right\rangle}
\DeclareMathOperator{\sgn}{sgn}
\numberwithin{equation}{section}
\allowdisplaybreaks[2]
\setlist[enumerate]{itemsep=3pt,topsep=5pt}
\setlist[itemize]{itemsep=3pt,topsep=5pt}

\title{The cotype--cotype conjecture\\under the approximation property}
\author{OpenAI}
\date{September 23, 2026}
\hypersetup{
  pdftitle={The cotype-cotype conjecture under the approximation property},
  pdfauthor={OpenAI},
  pdfsubject={Finite cotype, K-convexity, approximation, and Grothendieck pairs},
  pdfkeywords={Banach spaces, cotype, K-convexity, approximation property}
}
\begin{document}
\maketitle
\begin{abstract}
We prove the cotype--cotype conjecture under the ordinary approximation
property. A nonzero real Banach space with this property is \(K\)-convex
if and only if both the space and its dual have finite Rademacher cotype,
possibly with different exponents.
\end{abstract}

\section{The characterization}\label{sec:introduction}

The cotype--cotype problem asks whether finite Rademacher cotype
of a Banach space and its dual forces \(K\)-convexity. Without an
approximation hypothesis the assertion is false. We prove the characterization under the
ordinary approximation property.

All spaces in this paper are real, and \(X^*\) denotes the continuous
dual of \(X\). Recall that \(X\) has the \emph{approximation property}
(AP) if, for every compact set \(M\subset X\) and every \(\delta>0\),
there is a bounded finite-rank linear map \(S:X\to X\) such that
\[
\sup_{x\in M}\|Sx-x\|<\delta.
\]
The norms of these maps need not have a common bound. The
\emph{bounded approximation property} (BAP) requires, in addition,
a constant \(\lambda<\infty\) such that all the maps needed for these
approximations can be chosen with \(\|S\|\le\lambda\).

For \(2\le q<\infty\), let \(C_q(Y)\) be the least constant such that
\[
\left(\sum_{i=1}^s\|y_i\|^q\right)^{1/q}
\le C_q(Y)
\left(\E_\eps\left\|\sum_{i=1}^s\eps_i y_i\right\|^2\right)^{1/2}
\]
for every finite family in \(Y\), where the \(\eps_i\) are independent
uniform signs. The space \(Y\) has \emph{finite Rademacher cotype} if
\(C_q(Y)<\infty\) for some finite \(q\).

The \emph{Rademacher projection} on \(X\)-valued functions on a sign
cube retains their first-order coefficients:
\[
(\Rad_s f)(\eps)=\sum_{i=1}^s\eps_i\,\E_\eta[\eta_i f(\eta)].
\]
Write
\[
K(X)=\sup_{s\ge1}
\|\Rad_s:L_2(\{-1,1\}^s;X)\to L_2(\{-1,1\}^s;X)\|.
\]
Here \(L_2\) uses uniform probability measure. By definition, \(X\) is
\emph{\(K\)-convex} when \(K(X)<\infty\).

\begin{theorem}\label{thm:main}
Let \(X\) be a nonzero real Banach space with the ordinary approximation
property. Then \(X\) is \(K\)-convex if and only if
\[
C_q(X)<\infty
\quad\hbox{and}\quad
C_r(X^*)<\infty
\qquad\text{for some }q,r\in[2,\infty).
\]
The exponents \(q\) and \(r\) may differ. The forward implication holds
for every real Banach space, without an approximation hypothesis.
\end{theorem}

A space \(X\) has \emph{Rademacher type \(p\)}, \(1<p\le2\), if for some \(D<\infty\),
\[
\left(\E_\eps\left\|\sum_i\eps_i x_i\right\|^2\right)^{1/2}
\le D\left(\sum_i\|x_i\|^p\right)^{1/p}
\]
for all finite families. We say that \(X\) contains the \(\ell_1^s\)'s
uniformly if the isomorphism distortions of such copies inside \(X\)
are bounded independently of \(s\).

The classical theorems of Maurey--Pisier and Pisier
\cite{MP76,Pis82K} identify \(K\)-convexity with nontrivial type:
\begin{equation}\label{eq:classical-characterization}
\begin{split}
K(X)<\infty
&\ \Longleftrightarrow\
X\text{ has Rademacher type }p\text{ for some }p>1\\
&\ \Longleftrightarrow\
X\text{ does not contain the }\ell_1^s\text{'s uniformly}.
\end{split}
\end{equation}
These equivalences require no approximation hypothesis; precise
formulations are \cite[Theorem~4.40(i)]{Pis11IHP} and
\cite[Theorem~1.5]{Pis80Kseminar}.
The finite-cube definition above agrees with the infinite-product
definition in \cite[Definition~1.1]{Pis80Kseminar}, by density of
functions depending on finitely many coordinates.

The characterization is uniform when the cotype data are fixed. Its
finite-dimensional form also supplies a class of convex bodies for which
ordinary metric-entropy duality has dimension-free constants.
For an origin-symmetric convex body \(B\subset\R^n\), write \(X_B\)
for \(\R^n\) with unit ball \(B\), and \(B^\circ\) for its polar.
Here a convex body is compact with nonempty interior, and \(N(S,T)\)
is the least number of translates of \(T\) needed to cover \(S\).

\begin{corollary}[Uniform cotype bounds and entropy duality]
\label{cor:uniform-cotype-entropy}
Fix \(q,r\in[2,\infty)\) and \(C,D\in[1,\infty)\). There is a
finite constant \(\kappa=\kappa(q,r,C,D)\) such that every nonzero
real Banach space \(X\) with AP and
\[
C_q(X)\le C,\qquad C_r(X^*)\le D
\]
satisfies \(K(X)\le\kappa\).
Consequently there are \(a,b\ge1\), depending only on \(q,r,C,D\),
such that for every \(n\ge1\), all origin-symmetric convex bodies
\(K,L\subset\R^n\), and every \(t>0\),
\begin{equation}\label{eq:cotype-entropy-duality}
\frac1b\log N(L^\circ,atK^\circ)
\le \log N(K,tL)
\le b\log N(L^\circ,a^{-1}tK^\circ),
\end{equation}
provided either \(X_K\) or \(X_L\) satisfies the displayed primal
and dual cotype bounds.
\end{corollary}

The constants are existential; no bound on the AP approximants is
required. The entropy conclusion uses the bounded-\(K\) duality theorem
of Artstein--Milman--Szarek--Tomczak-Jaegermann
\cite[Theorem~5]{AMSTJ2004}, in the covering-number form of
\cite[Theorem~1.1 and Section~2]{Milman2007}.
We prove the corollary at the end of Section~\ref{sec:approximation}.

\subsection{History and scope}

Type and cotype express the geometry of a Banach space through the
size of random signed sums. Their relation to uniform finite-dimensional
subspaces was developed by Maurey and Pisier \cite{MP76}.
Elementary duality passes a type estimate to a cotype estimate on the
dual, but recovering type from cotype requires further information.
Maurey and Pisier introduced \(K\)-convexity as the setting for this
duality question; Pisier's holomorphic-semigroup theorem subsequently
identified it with nontrivial type
\cite{MP76,Pis80Kseminar,Pis82K}.
The question here asks whether finite cotype on both sides, together
with approximation, already puts the space in that setting.

Approximation assumptions have played a substantial role in these questions.
Enflo's counterexample showed that a separable reflexive Banach space can
fail AP \cite{Enf73}. A Schauder basis gives uniformly bounded finite-rank
approximations, but the converse implications in the hierarchy
\[
\text{Schauder basis}\ \Longrightarrow\ \text{BAP}\ \Longrightarrow\ \text{AP}
\]
fail: Figiel and Johnson separated AP from BAP \cite[p.~199]{FJ73}, and
Szarek constructed a superreflexive space with BAP but no Schauder basis
\cite[Theorem~1.1]{Sza87}.
Thus an argument using ordinary AP must accommodate approximants whose
norms have no common bound.

For cotype two on both sides, Pisier established the stronger conclusion
that an AP space is isomorphic to a Hilbert space
\cite[Corollaries~2.9 and~2.16]{Pis80}.
His subsequent work proved \(K\)-convexity under a restriction on the
cotype exponents, first with a basis or a dense increasing family of
uniformly complemented finite-dimensional subspaces, and then with BAP
\cite[Theorem~1, Remarks~2 and~9]{Pis81}.
He asked whether the exponent restriction could be removed
\cite[Remark~2(iii)]{Pis81}.
Szarek and Tomczak-Jaegermann discuss this question as the
\emph{cotype--cotype conjecture} in connection with uniformly complemented
finite-dimensional subspaces \cite{STJ09}.
Gupta, Misra, and Ray state the BAP formulation in
\cite[Conjecture~1.5]{GMR25}.
Theorem~\ref{thm:main} proves the characterization for all finite cotype
exponents under ordinary AP.

The assertion with no approximation hypothesis is false.
Pisier's construction embeds every cotype-two space into a space whose
primal and dual both have cotype two \cite[Theorem~3.2]{Pis83}.
Starting with \(\ell_1\) therefore gives a space which is not
\(K\)-convex; these examples fail AP \cite[Remark~3.6]{Pis83}.

Another application concerns a uniform tensor estimate. Given a sequence
\(\mathcal E=(E_n)\) with \(\dim E_n=n\), one asks whether every map
\(A:E_m\to E_n^*\), tensored with the identity of \(F\), is bounded
from the injective tensor norm to the projective tensor norm by
one constant times \(\|A\|\), independently of \(m,n,A\).
This is the Grothendieck-pair condition, stated formally in
Section~\ref{sec:grothendieck}. Gupta--Misra--Ray prove that it forces
finite cotype of \(F\) and \(F^*\), and record Pisier's conjecture
that for infinite-dimensional BAP spaces \(F\) the estimate then also
holds with \(\ell_2\) in place of \(F\)
\cite[Definition~1.2, Theorem~1.3 and Conjecture~1.4]{GMR25}.
Their deduction through complemented Euclidean subspaces
\cite[Corollary~4.8]{GMR25}, combined with Theorem~\ref{thm:main},
gives this conclusion under ordinary AP.
This uniform operator estimate differs from Grothendieck's original
injective/projective tensor-product conjecture, disproved in
\cite{Pis83}, and from the GT-pairs of \cite[Section~6]{Pis12}.

Rademacher and Walsh expansions connect probabilistic moment estimates
with Banach-space geometry. Kahane's inequality compares moments of
vector-valued Rademacher sums \cite[Theorem~8.1]{Pis11IHP};
Borell extended moment comparison to Banach-valued Walsh polynomials
of bounded degree \cite{Borell79}. The moment lifts below combine the
Rademacher comparison with quotient duality and Hahn--Banach. The scalar coefficient
lemma uses an additional feature of Fourier indices: their rank over
\(\F_2\), which determines the number of independent Walsh factors.

\subsection{The mechanism of the proof}

The proof has two parts. First, we establish a uniform decay estimate
for the normalized Walsh transforms of finite-rank operators
\(T:G\to F\): apply \(T\) to a vector family and then multiply by the
matrix \((2^{-n}(-1)^{u\cdot v})_{u,v\in\F_2^n}\), with the dot product
computed modulo two and probability-normalized \(L_2\) norms.
Cotype estimates for \(G^*\) and \(F\) at one sufficiently large, fixed
tuple length suffice. The decay estimate, relative to \(\|T\|\), is independent of the
operator rank and of the spaces; Theorem~\ref{thm:uniform-decay} states the
precise result.

A fixed finite-rank operator already has a Walsh bound obtained by
factoring through a Hilbert space, but its constant depends on the
operator. To remove that dependence, we show that a slowly decreasing
Walsh norm forces another substantial norm at nearly twice the
dimension. A moment lift represents a tuple \((x_i)\) by a bounded
sign-indexed family \((P_\eps)\) with
\(\E_\eps\eps_iP_\eps=ax_i\), where \(a>0\) is controlled by cotype.
Such lifts and finite-group averaging turn the slow decrease into
bounded scalar functions with prescribed row and column moments.
Their Fourier expansion is indexed by binary matrices and must contain
an index of rank at least two over \(\F_2\). Such an index produces
two independent Walsh factors.
Iteration would then contradict the Hilbert-space bound for the same
fixed operator. This proves uniform decay.

Second, we transfer that decay to the identity on \(X\) using only AP.
For each finite-dimensional \(E\subset X\), we construct a compact map
\(J:G\to X\) of uniformly bounded norm and a contraction \(\iota:E\to G\)
with \(J\iota x=x\). The dual of \(G\) has a controlled moment estimate
(Lemma~\ref{lem:compact-domain}). The construction uses moment lifts in
\(X\), rather than assuming that the quotient \(E^*\) inherits the
cotype of \(X^*\). AP gives finite-rank composites \(S_jJ\) converging
to \(J\) in operator norm. Uniform Walsh decay passes to \(J\), and
then to every finite Walsh input in \(X\). Uniform copies of
\(\ell_1^s\) would contradict that decay.

Section~\ref{sec:preliminaries} introduces Walsh norms and the moment lifts.
Section~\ref{sec:coefficient} proves the scalar coefficient lemma, and
Section~\ref{sec:walsh} turns it into uniform operator decay.
Section~\ref{sec:approximation} constructs the compact domains and proves
Theorem~\ref{thm:main} and Corollary~\ref{cor:uniform-cotype-entropy}.
Section~\ref{sec:grothendieck} gives the tensor-norm consequence.

\section{Walsh transforms and moment lifts}\label{sec:preliminaries}

We first introduce the operator norms that measure Walsh decay, then derive
the moment lifts needed to compare those norms at different dimensions.
Throughout, Banach spaces are real. All unspecified averages on finite
sets are uniform, and distinct averaging variables are independent.
The notation \(B_Y\) means the closed unit ball of \(Y\).
For a finite probability space \(\Omega\), the space \(L_p(\Omega;Y)\)
has its probability-normalized norm. By contrast,
\(\ell_2^m(Y)\) has norm \((\sum_{i=1}^m\|y_i\|^2)^{1/2}\).
Duality on finite probability spaces uses the expectation pairing;
in particular,
\[
\begin{gathered}
L_2(\Omega;Y)^*=L_2(\Omega;Y^*),\\
L_1(\Omega;Y)^*=L_\infty(\Omega;Y^*),\qquad
L_\infty(\Omega;Y)^*=L_1(\Omega;Y^*).
\end{gathered}
\]

\subsection{Walsh norms}

Let \(V_n=\F_2^n\), with its standard dot pairing, and put
\(\chi(u,v)=(-1)^{u\cdot v}\). For \(T:G\to F\), define
\[
(W_n^Tx)(v)=\E_{u\in V_n}\chi(u,v)Tx(u),\qquad
w_n(T)=\|W_n^T:L_2(V_n;G)\to L_2(V_n;F)\|.
\]
We allow \(n=0\), so \(w_0(T)=\|T\|\).
Finite-sum duality gives the equivalent formula, with
\(x_u\in G\) and \(y_v\in F^*\),
\begin{equation}\label{eq:walsh-bilinear}
w_n(T)=\sup\left|
\E_{u,v\in V_n}\chi(u,v)y_v(Tx_u)\right|,
\qquad
\E_u\|x_u\|^2\le1,\quad \E_v\|y_v\|^2\le1.
\end{equation}

\begin{lemma}[Basic Walsh estimates]\label{lem:walsh-basic}
For every bounded \(T:G\to F\):
\begin{enumerate}[label=\textup{(\roman*)}]
\item \(2^{-n/2}\|T\|\le w_n(T)\le\|T\|\), and \(w_n(T)\) is
nonincreasing in \(n\).
\item If \(T_\Omega:L_2(\Omega;G)\to L_2(\Omega;F)\) acts pointwise
as \(T\), where \(\Omega\) is finite, then
\(w_n(T_\Omega)=w_n(T)\).
\item If \(T\) has finite rank, there is \(D_T<\infty\) such that
\begin{equation}\label{eq:rank-dependent-decay}
w_n(T)\le D_T2^{-n/2}\qquad(n\ge0).
\end{equation}
\end{enumerate}
Moreover, for \(A,B:G\to F\),
\begin{equation}\label{eq:walsh-continuity}
|w_n(A)-w_n(B)|\le\|A-B\|.
\end{equation}
\end{lemma}
\begin{proof}
The upper bound follows from the triangle inequality and
Cauchy--Schwarz; the lower bound follows by taking input supported
at one \(u\), scaled by \(2^{n/2}\). On \(V_{n+j}=V_n\oplus V_j\),
first apply \(W_n^T\) in the first coordinate and then the normalized
Walsh transform with the identity of \(F\) in the second coordinate.
The latter has norm at most \(1\). Averaging the squared estimates
in the unused coordinates proves monotonicity.
Interchanging finite \(L_2\) coordinates proves (ii), with constant
functions on \(\Omega\) giving the reverse inequality.
For (iii), factor a finite-rank \(T\) as \(G\to H\to F\), with \(H\)
a finite-dimensional Hilbert space. Orthogonality of the Walsh
matrix gives \(\|W_n^{\Id_H}\|=2^{-n/2}\), proving (iii).
Finally, \(W_n^A-W_n^B=W_n^{A-B}\), so the upper bound in (i)
and the reverse triangle inequality prove \eqref{eq:walsh-continuity}.
\end{proof}

\subsection{Fixed-length cotype estimates}

Define \(c_m(Y)\) to be the least constant in
\begin{equation}\label{eq:fixed-cotype}
\left(\sum_{i=1}^m\|y_i\|^2\right)^{1/2}
\le c_m(Y)\left(\E_\eps\left\|\sum_{i=1}^m\eps_i y_i\right\|^2\right)^{1/2}.
\end{equation}
Thus, if \(Y\) has cotype \(q\) with constant \(C_q(Y)\), then
\begin{equation}\label{eq:fixed-from-cotype}
c_m(Y)\le C_q(Y)m^{1/2-1/q}.
\end{equation}
Squaring \eqref{eq:fixed-cotype} and averaging proves
\begin{equation}\label{eq:cotype-amplification}
c_m(L_2(\Omega;Y))\le c_m(Y).
\end{equation}
The same observation applies to duals of these finite \(L_2\) sums.

The following estimate is the \(L_1\)-to-\(L_2\) case of Kahane's
inequality; see \cite[Theorem~8.1]{Pis11IHP}.
We include a stopping-time proof with an explicit, sufficient constant
so that the subsequent moment lifts are self-contained.

\begin{lemma}[Moment comparison]\label{lem:kahane}
For every finite Rademacher sum \(S=\sum_{i=1}^m\eps_i z_i\) in a Banach
space,
\begin{equation}\label{eq:kahane}
(\E\|S\|^2)^{1/2}\le L\,\E\|S\|,\qquad L=24.
\end{equation}
\end{lemma}
\begin{proof}
Put \(M=\E\|S\|\) and \(b=\max_i\|z_i\|\).
Since \(z_i=\E(\eps_iS)\), we have \(b\le M\).
The case \(M=0\) is immediate.

Let \(S_j=\sum_{i\le j}\eps_i z_i\), and stop at the first \(j\) with
\(\|S_j\|>u\). Conditional on a stopped prefix \(z\), the remaining
sum is symmetric; at least one of \(\|z+h\|\) and \(\|z-h\|\)
exceeds \(u\). Hence
\[
\mathbb P(\text{a stop occurs})\le2\mathbb P(\|S\|>u).
\]
For any fixed \(j\), symmetry of the independent prefix similarly gives
\[
\mathbb P\left(\left\|\sum_{i>j}\eps_i z_i\right\|>v\right)
\le2\mathbb P(\|S\|>v).
\]
At the first stop, \(\|S_j\|\le u+b\). Decomposing according to its
time and using independence of the remaining sum therefore yields
\[
\mathbb P(\|S\|>u+b+v)
\le4\mathbb P(\|S\|>u)\mathbb P(\|S\|>v).
\]
Take \(u=8M\). Markov's inequality and \(b\le M\) imply
\[
\mathbb P(\|S\|>9M+v)\le\tfrac12\mathbb P(\|S\|>v).
\]
Thus \(\mathbb P(\|S\|>9jM)\le2^{-j}\) for \(j\ge0\), and integration
over consecutive intervals of length \(9M\) gives
\[
\E\|S\|^2
\le81M^2\sum_{j\ge0}(2j+1)2^{-j}
=486M^2<24^2M^2.
\]
\end{proof}

We will repeatedly use a simple consequence of quotient duality.
If a bounded surjection \(U:A\to B\) satisfies
\(\|U^*b^*\|\ge D^{-1}\|b^*\|\), then
\begin{equation}\label{eq:quotient-lift}
\inf_{Ua=b}\|a\|\le D\|b\|.
\end{equation}
Indeed, \(A/\ker U\) is isomorphic to \(B\), and its quotient norm at
\(b\) is
\[
\sup_{\|U^*b^*\|\le1}|\langle b,b^*\rangle|.
\]
In each application below surjectivity has an explicit bounded right
inverse, so the quotient identification is immediate.

\begin{lemma}[Bounded moment lifts]\label{lem:lifts}
Suppose
\[
c_m(G^*)\le D_g,\qquad c_m(F)\le D_f,
\qquad D_g,D_f\ge1.
\]
Set
\begin{equation}\label{eq:lift-parameters}
a=\frac1{2LD_g\sqrt m},\qquad
b=\frac1{2LD_f\sqrt m}.
\end{equation}
Then:
\begin{enumerate}[label=\textup{(\roman*)}]
\item For \(x_1,\ldots,x_m\in B_G\), there are \(P_\eps\in B_G\) such that
\[
\E_\eps\eps_iP_\eps=ax_i\quad(1\le i\le m).
\]
\item For \(y_1,\ldots,y_m\in B_{F^*}\), there are \(Z_\delta\in B_{F^*}\)
such that
\[
\E_\delta\delta_jZ_\delta=by_j\quad(1\le j\le m).
\]
\end{enumerate}
The same conclusions, with the same constants, hold after replacing
\(G,F\) by finite \(L_2\) sums of their copies.
\end{lemma}
\begin{proof}
The moment map
\[
U:L_\infty(\{-1,1\}^m;G)\longrightarrow\ell_2^m(G),
\qquad
UP=(\E_\eps\eps_iP_\eps)_{i=1}^m,
\]
is onto: a right inverse is \((x_i)\mapsto\sum_i\eps_i x_i\).
Its adjoint sends \((z_i^*)\) to \(\sum_i\eps_i z_i^*\).
By Lemma~\ref{lem:kahane} and \eqref{eq:fixed-cotype},
\[
\|U^*(z_i^*)\|_{L_1(G^*)}
\ge(LD_g)^{-1}\left(\sum_i\|z_i^*\|^2\right)^{1/2}.
\]
Equation \eqref{eq:quotient-lift} bounds the infimal lift norm of
\((ax_i)\) by \(LD_ga\sqrt m=1/2\). A lift of norm at most \(1\)
therefore exists, proving (i).

For (ii), the map
\[
V:\ell_2^m(F)\longrightarrow L_1(\{-1,1\}^m;F),
\qquad
V(w_j)=\sum_j\delta_jw_j,
\]
satisfies
\(\|Vw\|_{L_1}\ge(LD_f)^{-1}\|w\|_{\ell_2^m(F)}\).
The functional
\[
Vw\longmapsto b\sum_j y_j(w_j)
\]
has norm at most \(b\sqrt m\,LD_f=1/2\) on its range.
Hahn--Banach extends it to \(L_1(F)\) with the same norm.
Representing the extension in \(L_\infty(F^*)\) gives \(Z\) and the
stated moments. The amplification assertion follows from
\eqref{eq:cotype-amplification} and finite-sum duality.
\end{proof}

\section{A scalar Fourier-coefficient lemma}\label{sec:coefficient}

For a function on a sign cube, write
\[
\Odd h(z)=\frac{h(z)-h(-z)}2,
\]
where \(-z\) negates every coordinate. We call \(h\) odd if \(\Odd h=h\).
Fourier coefficients always use uniform measure. For an \(m\times m\)
sign matrix \(S\) and a binary matrix \(P\in\{0,1\}^{m\times m}\), put
\[
S^P=\prod_{i,j}S_{ij}^{P_{ij}},\qquad
\widehat h(P)=\E_S h(S)S^P.
\]
For vector cubes we use binary indices \(p\in\{0,1\}^m\);
\(|p|\) denotes the number of ones and \(e_i\) the \(i\)-th unit vector.
The notation \(\sigma S\) multiplies row \(i\) by \(\sigma_i\), and
\(S\tau\) multiplies column \(j\) by \(\tau_j\).
We denote row \(i\) and column \(j\) by \(S_{i,:}\) and \(S_{:,j}\), respectively.
All ranks in this section are over \(\F_2\).

The point of the binary rank is the identity
\[
\prod_{i,j}\chi(u_i,v_j)^{P_{ij}}
=(-1)^{\sum_{i,j}P_{ij}u_i\cdot v_j},
\qquad u_i,v_j\in\F_2^n.
\]
An index of rank \(k\) yields \(k\) independent bilinear terms after
binary changes of variables. In Section~\ref{sec:walsh}, this will turn a
nonzero coefficient into a Walsh pairing in \(k\) times as many
coordinates. We therefore need a coefficient whose index has rank
at least two.

\par\medskip\noindent\begin{minipage}{\textwidth}
For example, take \(S_{ij}=\chi(u_i,v_j)\) for \(1\le i,j\le3\).
The two indices with an odd number of ones
\[
\begin{gathered}
P_1=\underbrace{\begin{pmatrix}
1&1&1\\1&1&1\\1&1&1
\end{pmatrix}}_{\text{binary rank }1},
\qquad S^{P_1}=\chi(u_1+u_2+u_3,v_1+v_2+v_3),\\[6pt]
P_2=\underbrace{\begin{pmatrix}
1&1&0\\0&1&0\\0&0&0
\end{pmatrix}}_{\text{binary rank }2},
\qquad S^{P_2}=\chi(u_1,v_1+v_2)\chi(u_2,v_2)
\end{gathered}
\]
yield one and two Walsh factors, respectively: \(v_1+v_2\) and
\(v_2\) vary independently. \end{minipage}\par\medskip
Rank-one indices have rectangular support.
The next lemma rules out an expansion supported only on odd rectangles
when both families of prescribed moments are large enough.

\begin{lemma}[A coefficient of rank at least two]\label{lem:coefficient}
Let \(M_0=4\), and suppose that
\[
4\le d=\lceil2\log_2m\rceil<m.
\]
Let \(a,b>0\), put \(A=ma\), \(B=mb\), and assume
\begin{equation}\label{eq:large-parameters}
\min(A,B)\ge20M_0d(d+1)^2.
\end{equation}
There is \(c_*>0\), depending only on \(m,a,b\), with the following
property. Suppose that
\[
\Psi:\{-1,1\}^{m\times m}\to\R,\qquad
f,g:\{-1,1\}^m\to\R
\]
are odd, have sup norms at most \(M_0\), and satisfy
\begin{equation}\label{eq:scalar-first}
\widehat f(e_i)=a,\qquad \widehat g(e_j)=b
\end{equation}
together with the projection identities
\begin{equation}\label{eq:scalar-projections}
\begin{split}
\E_\sigma\sigma_i\Psi(\sigma S)&=a\,g(S_{i,:}),\\
\E_\tau\tau_j\Psi(S\tau)&=b\,f(S_{:,j})
\end{split}
\end{equation}
for all \(i,j,S\). Then
\[
|\widehat\Psi(P)|\ge c_*
\]
for some binary matrix \(P\) of rank at least two.
\end{lemma}

\begin{proof}
We first reduce the uniform lower bound to a nonvanishing assertion.
If every index in the Fourier support of \(\Psi\) had rank at most one, a damped polynomial would encode
their rectangular supports. Its prescribed mixed derivative will
contradict a bound obtained by splitting it into two parts, each of
low degree in one variable.

The admissible triples \((\Psi,f,g)\) form a compact set in a
finite-dimensional space. The maximum of
\(|\widehat\Psi(P)|\) over \(\rank P\ge2\) is continuous.
It therefore suffices to show that this maximum never vanishes.
If the admissible set is empty, any positive \(c_*\) will do.

Suppose, to the contrary, that the Fourier support of \(\Psi\)
has rank at most one. Oddness eliminates indices with an even
number of ones. Every remaining index has a unique form
\(P=ph^t\) with \(p,h\ne0\), and both \(k=|p|\) and \(l=|h|\) are odd.
Let
\[
C_{kl}=\sum_{\substack{|p|=k\\|h|=l}}\widehat\Psi(ph^t).
\]
Parseval's identity and Cauchy--Schwarz give
\begin{equation}\label{eq:Ckl-bound}
|C_{kl}|\le M_0\sqrt{\binom mk\binom ml}.
\end{equation}
To suppress rectangles that are large in both directions, consider
the damped polynomial
\[
H(s,u)=\sum_{\substack{1\le k,l\le m\\ k,l\ {\rm odd}}}
C_{kl}s^ku^l2^{-kl}.
\]
It has sup norm at most \(M_0\) on \([-1,1]^2\).
To see this, take independent signs \(D_i,E_j,R_{ij}\) with
means \(s,u,1/2\), respectively, and set \(S_{ij}=D_iE_jR_{ij}\).
For an odd rectangle \(ph^t\),
\[
\E S^{ph^t}=s^{|p|}u^{|h|}2^{-|p||h|},
\]
so \(H(s,u)=\E\Psi(S)\).

Delete from \(H\) the terms with both \(k>d\) and \(l>d\), obtaining
\(H_0\). If \(k,l>d\), then
\[
(k+l)\log_2m\le kl,\qquad kl/2\ge k+l.
\]
Using \eqref{eq:Ckl-bound}, the deleted part is bounded in sup norm by
\[
M_0\sum_{k,l>d}m^{(k+l)/2}2^{-kl}
\le M_0\sum_{k,l>d}2^{-kl/2}
\le M_0\sum_{k,l>d}2^{-k-l}\le M_0.
\]
Consequently \(\|H_0\|_\infty\le2M_0\).

For an odd rectangle \(ph^t\), the vectors of row and column parities
are \(p\) and \(h\). Comparing Fourier coefficients in
\eqref{eq:scalar-projections} therefore gives
\[
C_{1l}=ma\sum_{|h|=l}\widehat g(h),\qquad
C_{k1}=mb\sum_{|p|=k}\widehat f(p).
\]
The deletion defining \(H_0\) does not affect terms with \(k=1\) or
\(l=1\). Hence
\begin{equation}\label{eq:H-derivatives}
\begin{split}
\partial_sH_0(0,u)&=A\sum_h\widehat g(h)(u/2)^{|h|},\\
\partial_uH_0(s,0)&=B\sum_p\widehat f(p)(s/2)^{|p|},\\
\partial_s\partial_uH_0(0,0)&=AB/2.
\end{split}
\end{equation}
Independent signs of common mean \(u/2\), or \(s/2\), show that
the first two expressions have sup norms at most \(AM_0\) and \(BM_0\).

Every retained monomial has low degree in at least one variable.
We separate these two possibilities with controlled sup norms.
Expand \(H_0\) in the Chebyshev basis \(T_j(s)\), defined by
\(T_j(\cos\theta)=\cos(j\theta)\). Let \(H_1\) retain the terms
with \(j\le d\), and set \(H_2=H_0-H_1\).
Then
\begin{equation}\label{eq:separate-degrees}
\deg_sH_1\le d,\qquad \deg_uH_2\le d.
\end{equation}
Indeed, every monomial of \(H_0\) has \(s\)-degree at most \(d\)
or \(u\)-degree at most \(d\), and truncation in the Chebyshev
basis annihilates every monomial of the former kind in \(H_2\).

Each Chebyshev coefficient of a polynomial on \([-1,1]\) is bounded
by twice its sup norm, by cosine orthogonality. Since \(|T_j|\le1\),
\[
\|H_1\|_\infty\le4M_0(d+1),\qquad
\|H_2\|_\infty\le6M_0(d+1).
\]
Also, \(|T_j'(0)|\le j\), so every polynomial \(Q\) of degree at most
\(d\) satisfies
\begin{equation}\label{eq:derivative-bound}
|Q'(0)|\le d(d+1)\|Q\|_\infty.
\end{equation}
First apply \eqref{eq:derivative-bound} in \(s\) to \(H_1\), and
in \(u\) to \(H_2\). Then \eqref{eq:H-derivatives} implies
\[
\begin{split}
\|\partial_sH_2(0,\cdot)\|_\infty
&\le AM_0+4M_0d(d+1)^2,\\
\|\partial_uH_1(\cdot,0)\|_\infty
&\le BM_0+6M_0d(d+1)^2.
\end{split}
\]
By \eqref{eq:separate-degrees}, these are polynomials of degree
at most \(d\) in their displayed variables. Applying
\eqref{eq:derivative-bound} once more and adding yields
\begin{equation}\label{eq:mixed-upper}
|\partial_s\partial_uH_0(0,0)|
\le M_0d(d+1)\bigl(A+B+10d(d+1)^2\bigr).
\end{equation}
On the other hand, \eqref{eq:large-parameters} gives
\[
\frac{AB}{2}\ge5M_0d(d+1)^2(A+B),\qquad
10d(d+1)^2\le\frac{A+B}{4M_0}.
\]
Since \(5(d+1)>1+1/(4M_0)\), these inequalities contradict
\eqref{eq:H-derivatives} and \eqref{eq:mixed-upper}.
\end{proof}

\section{Uniform Walsh decay for finite-rank operators}\label{sec:walsh}

\subsection{Dimension growth}

We now show that a slow decrease of Walsh norms forces a large Walsh
norm at almost twice the original index. The proof uses the scalar
coefficient lemma to find several independent Walsh factors inside a
single pairing. To make that lemma applicable, we first arrange that
our pairings depend only on their matrices of signs. The initial Walsh
test is indexed by vectors, whereas the scalar lemma concerns arbitrary
sign matrices; the frame construction below removes the extra dependence
on the vectors.

An ordered linearly independent tuple in \(V_n\) will be called a
\emph{frame}. We identify covectors with vectors through the dot
pairing. Independence is required separately on the two sides.

\begin{lemma}[Pairs of frames]\label{lem:frames}
The action
\[
(U,V)\longmapsto(HU,H^{-t}V),\qquad H\in\GL(n,\F_2),
\]
is transitive on pairs of frames of fixed lengths with a prescribed
matrix of pairings \((u_i\cdot v_j)\).
Every binary \(r\times s\) matrix is such a pairing matrix if
\(n\ge r+s\).
\end{lemma}
\begin{proof}
Consider two pairs \((U,V)\), \((U',V')\) with the same pairing
matrix. The covector frames define surjections
\(\pi,\pi':V_n\to\F_2^s\).
The isomorphism \(u_i\mapsto u_i'\) between the primal spans respects
evaluation, and hence maps
\(\Span(U)\cap\ker\pi\) onto \(\Span(U')\cap\ker\pi'\).
Extend this restriction to an isomorphism between the two kernels.
Together with the map on the primal spans, it defines an isomorphism
\[
\Span(U)+\ker\pi\longrightarrow\Span(U')+\ker\pi'
\]
satisfying \(\pi'H=\pi\).
Complete it to an invertible map of \(V_n\) by choosing matching
lifts of a basis complementary to the common evaluation image
\(\pi(\Span(U))=\pi'(\Span(U'))\).
Then \(HU=U'\) and \(\pi'H=\pi\), the latter being equivalent to
\(H^{-t}V=V'\).
For feasibility in dimension \(r+s\), given a matrix \(D\), take
\(u_i=(e_i,0)\) and \(v_j=(D_{:,j},e_j)\).
\end{proof}

\begin{lemma}[Growth from a slow decrease]\label{lem:growth}
Fix an integer \(m\ge1\) and constants \(D_g,D_f\ge1\).
Suppose that \(m\) and the parameters \(a,b\)
in \eqref{eq:lift-parameters} satisfy the hypotheses of
Lemma~\ref{lem:coefficient}. There is \(c_1\in(0,1/2]\), depending only
on these parameters, such that the following holds.
If
\[
c_m(G^*)\le D_g,\qquad c_m(F)\le D_f,\qquad 0\ne T:G\to F,
\]
and \(n>4m\) satisfies
\begin{equation}\label{eq:slow-decrease}
w_{n-2m}(T)\le2w_n(T),
\end{equation}
then
\begin{equation}\label{eq:growth}
w_{2(n-2m)}(T)\ge c_1w_n(T).
\end{equation}
The operator \(T\) need not have finite rank.
\end{lemma}

\begin{proof}
Put \(n'=n-2m\). We will turn a Walsh test at index \(n\) into
bounded scalar functions satisfying Lemma~\ref{lem:coefficient}. A
coefficient of binary rank \(k\ge2\) will then give a Walsh test
at index \(kn'\).

\medskip
\textit{Step 1: realize the Walsh signs exactly and equivariantly.}
Choose families \(p_u\in G\), \(q_v\in F^*\) admissible in
\eqref{eq:walsh-bilinear}, changing a sign if necessary, such that
\[
\gamma=\E_{u,v}\chi(u,v)q_v(Tp_u)\ge w_n(T)/2>0.
\]
Amplify \(T\) over \(V_n\times V_n\), and define
\[
x'_u(\alpha,\beta)=(-1)^{u\cdot\beta+\alpha\cdot\beta}p_{u+\alpha},
\qquad
y'_v(\alpha,\beta)=(-1)^{\alpha\cdot v}q_{v+\beta}.
\]
Uniform translation shows that every \(x'_u\) and \(y'_v\) has
norm at most one in the corresponding finite \(L_2\) space.
Writing \(T'\) for the amplified operator, we obtain
\[
y'_v(T'x'_u)=\gamma\chi(u,v).
\]
Indeed, the pairing includes expectation over \(\alpha,\beta\),
and its phase is
\[
u\cdot\beta+\alpha\cdot\beta+\alpha\cdot v
=(u+\alpha)\cdot(v+\beta)+u\cdot v.
\]

Amplify once more over \(\mathscr G=\GL(n,\F_2)\), obtaining
\(\widetilde T:\widetilde G\to\widetilde F\), and set
\[
x_u(R)=x'_{Ru},\qquad y_v(R)=y'_{R^{-t}v}.
\]
For the bilinear form \(B(x,y)=y(\widetilde T x)\), this gives
\begin{equation}\label{eq:exact-pairing}
x_u\in B_{\widetilde G},\quad y_v\in B_{\widetilde F^*},
\qquad B(x_u,y_v)=\gamma\chi(u,v).
\end{equation}
Use the same right translations
\(\Lambda(H)z(R)=z(RH)\) on the domain and target dual.
They are isometries, preserve \(B\) under simultaneous translation,
and satisfy
\[
\Lambda(H)x_u=x_{Hu},\qquad
\Lambda(H)y_v=y_{H^{-t}v},\qquad
\Lambda(H_1)\Lambda(H_2)=\Lambda(H_1H_2).
\]
The orders follow directly from \((RH)^{-t}=R^{-t}H^{-t}\).
By Lemma~\ref{lem:walsh-basic},
\(w_j(\widetilde T)=w_j(T)\) for every \(j\).
Equation \eqref{eq:cotype-amplification} also preserves the bounds
\(D_g,D_f\), so the moment lifts remain available with the same
\(a,b\).

\medskip
\textit{Step 2: lift frames while preserving the symmetry.}
For every primal \(m\)-frame \(U=(u_i)_{i=1}^m\), choose
\(P_{U,\eps}\in B_{\widetilde G}\) with
\[
\E_\eps\eps_iP_{U,\eps}=ax_{u_i}.
\]
We can make these choices equivariant. For one reference frame,
take the lifts from Lemma~\ref{lem:lifts} and average each of them over
its stabilizer under \(\Lambda\). The prescribed moments are fixed
by this stabilizer, and the averaged lifts are stabilizer-invariant.
Transport them to any other frame by \(\Lambda(H)\).
Two transporters differ on the right by a stabilizer element, so
the result is well-defined and equivariant.

The same construction for covector frames, using the action
\(H^{-t}\), gives \(Z_{V,\delta}\in B_{\widetilde F^*}\) with
\[
\E_\delta\delta_jZ_{V,\delta}=by_{v_j}.
\]
The two systems of lifts can be chosen independently.
Their equivariance, invariance of \(B\), and Lemma~\ref{lem:frames}
show that the following pairings depend only on the displayed sign
patterns and lift labels:
\begin{equation}\label{eq:three-functions}
\begin{aligned}
\Phi(S,\eps,\delta)&=B(P_{U,\eps},Z_{V,\delta}),
& S_{ij}&=\chi(u_i,v_j),\\
\phi(s,\eps)&=B(P_{U,\eps},y_v),
& s_i&=\chi(u_i,v),\quad v\ne0,\\
\psi(t,\delta)&=B(x_u,Z_{V,\delta}),
& t_j&=\chi(u,v_j),\quad u\ne0.
\end{aligned}
\end{equation}
All patterns are feasible since \(n>4m\).
The moment identities and \eqref{eq:exact-pairing} give
\begin{equation}\label{eq:three-moments}
\begin{aligned}
\E_\eps\eps_i\Phi(S,\eps,\delta)&=a\psi(S_{i,:},\delta),
&
\E_\delta\delta_j\Phi(S,\eps,\delta)&=b\phi(S_{:,j},\eps),\\
\E_\eps\eps_i\phi(s,\eps)&=a\gamma s_i,
&
\E_\delta\delta_j\psi(t,\delta)&=b\gamma t_j.
\end{aligned}
\end{equation}

\medskip
\textit{Step 3: control the odd parts on a smaller Walsh cube.}
We need a bound that remains useful after division by \(\gamma\).
For every fixed choice of lift labels, we claim that
\begin{equation}\label{eq:odd-bounds}
\|\Odd_S\Phi(\cdot,\eps,\delta)\|_\infty,\quad
\|\Odd_s\phi(\cdot,\eps)\|_\infty,\quad
\|\Odd_t\psi(\cdot,\delta)\|_\infty
\le w_{n'}(T).
\end{equation}
Decompose \(V_n=V_{n'}\oplus\F_2^m\oplus\F_2^m\), and fix a sign
matrix \(S_0=((-1)^{D_{ij}})\).
For \(u,v\in V_{n'}\), take the primal frame
\[
U(u)_i=(u,e_i,0)
\]
and the two covector frames
\[
V^+(v)_j=(v,D_{:,j},e_j),\qquad
V^-(v)_j=(v,D_{:,j}+\mathbf1,e_j),
\]
where \(\mathbf1\in\F_2^m\) is the all-ones vector.
The tags \(e_i,e_j\) ensure independence for every \(u,v\).
Their sign matrices are \(\chi(u,v)S_0\) and
\(-\chi(u,v)S_0\), so
\[
B\left(P_{U(u),\eps},
\frac{Z_{V^+(v),\delta}-Z_{V^-(v),\delta}}2\right)
=\chi(u,v)(\Odd_S\Phi)(S_0,\eps,\delta).
\]
Both vector families lie in their respective unit balls, and each
depends only on its own Walsh index. Testing
\eqref{eq:walsh-bilinear} at \(n'\), and using amplification
invariance, proves the first bound.

For \(\phi\), retain \(U(u)\) and use the nonzero covectors
\((v,h,e_1)\), \((v,h+\mathbf1,e_1)\), where
\(s_i=(-1)^{h_i}\). Half the difference of the corresponding
\(y\)'s gives the same calculation.
For \(\psi\), use the covector frame \((v,0,e_j)\) and half the
difference of the primal vectors at
\((u,e_1,h)\), \((u,e_1,h+\mathbf1)\).
Their pairings are \(\chi(u,v)(-1)^{h_j}\) and its negative.
The \(e_1\) tags keep the singleton vectors nonzero.
This proves \eqref{eq:odd-bounds}.

\medskip
\textit{Step 4: find a Fourier coefficient of binary rank at least two.}
Define
\begin{equation}\label{eq:normalized-functions}
\begin{split}
f(s)&=\gamma^{-1}\Odd_s\E_\eps\phi(\eps s,\eps),\\
g(t)&=\gamma^{-1}\Odd_t\E_\delta\psi(\delta t,\delta),\\
\Psi(S)&=\gamma^{-1}\Odd_S\E_{\eps,\delta}
                  \Phi(\eps S\delta,\eps,\delta).
\end{split}
\end{equation}
These functions are odd, and the slow-decrease hypothesis gives
\[
\max(\|f\|_\infty,\|g\|_\infty,\|\Psi\|_\infty)
\le w_{n-2m}(T)/\gamma\le4.
\]
Changing variables \(s'=\eps s\) in the first coefficient and using
\eqref{eq:three-moments}, we obtain
\[
\widehat f(e_i)
=\gamma^{-1}\E_{s',\eps}s'_i\eps_i\phi(s',\eps)=a.
\]
The odd projection leaves this coefficient unchanged.
Likewise, \(\widehat g(e_j)=b\).

We next check the row and column projections required by
Lemma~\ref{lem:coefficient}. A first row projection is unchanged by
applying global odd projection, since
\[
\E_\sigma\sigma_i h(-\sigma S)
=-\E_\sigma\sigma_i h(\sigma S).
\]
Substitute \(\xi=\eps\sigma\); then \(\xi,\eps,\delta\) are still
independent uniform sign vectors. By \eqref{eq:three-moments},
\begin{align*}
\gamma\E_\sigma\sigma_i\Psi(\sigma S)
&=\E_{\xi,\eps,\delta}\xi_i\eps_i
                   \Phi(\xi S\delta,\eps,\delta)\\
&=a\,\E_{\xi,\delta}\xi_i
                  \psi(\xi_i(\delta S_{i,:}),\delta)\\
&=a\,\E_\delta\frac{\psi(\delta S_{i,:},\delta)
                         -\psi(-\delta S_{i,:},\delta)}2\\
&=a\gamma\,g(S_{i,:}).
\end{align*}
Interchanging rows and columns, and simultaneously
\((a,\eps,\phi)\) with \((b,\delta,\psi)\), gives
\(\E_\tau\tau_j\Psi(S\tau)=b f(S_{:,j})\).
Thus all hypotheses of Lemma~\ref{lem:coefficient} hold, with the fixed
parameters \(m,a,b\).

Choose a binary matrix \(P\) of rank \(k\ge2\) with
\(|\widehat\Psi(P)|\ge c_*\).
Since \(\Psi\) is odd and this coefficient is nonzero, the total
number of ones in \(P\) is odd. Its coefficient is therefore
unchanged when the odd projection is removed from
\eqref{eq:normalized-functions}.
Writing \(\Phi_{\eps,\delta}(S)=\Phi(S,\eps,\delta)\), row and column
sign multiplication gives the exact identity
\[
\gamma\widehat\Psi(P)
=\E_{\eps,\delta}
\left[\left(\prod_{i,j}(\eps_i\delta_j)^{P_{ij}}\right)
\widehat{\Phi_{\eps,\delta}}(P)\right].
\]
For some fixed \(\eps,\delta\), therefore,
\(h=\Phi_{\eps,\delta}\) satisfies
\begin{equation}\label{eq:h-coefficient}
|\widehat h(P)|\ge c_*\gamma.
\end{equation}
No oddness assertion about this individual function \(h\) is needed.

\medskip
\textit{Step 5: extract the independent Walsh factors.}
Keep \(P,\eps,\delta\) fixed. For arbitrary tuples
\((u'_i)_{i=1}^m\), \((v'_j)_{j=1}^m\) in \(V_{n'}\), set
\[
U_i=(u'_i,e_i,0),\qquad
V(D)_j=(v'_j,D_{:,j},e_j),
\]
where \(D\) ranges over binary \(m\times m\) matrices.
The tagged tuples are frames even if the original tuples have
repeated or zero entries.
The Fourier coefficient formula gives
\begin{equation}\label{eq:coefficient-extraction}
B\left(P_{U,\eps},
\E_D(-1)^{\langle P,D\rangle}Z_{V(D),\delta}\right)
=\widehat h(P)(-1)^{\sum_{i,j}P_{ij}u'_i\cdot v'_j},
\end{equation}
where \(\langle P,D\rangle=\sum_{i,j}P_{ij}D_{ij}\) is computed modulo
two. Indeed, translate the uniform matrix \(D\) by
\((u'_i\cdot v'_j)_{ij}\) in the defining Fourier average.
Both vectors in \eqref{eq:coefficient-extraction} have norm at most
one. The first depends only on the \(u'\)-tuple, and the second only
on the \(v'\)-tuple.

Since \(\rank_{\F_2}P=k\), choose binary \(m\times k\) matrices
\(I,J\) with \(I^tPJ=\Id_k\).
For \(\zeta,\eta\in(V_{n'})^k\), substitute
\[
u'_i=\sum_{\ell=1}^k I_{i\ell}\zeta_\ell,\qquad
v'_j=\sum_{\ell=1}^k J_{j\ell}\eta_\ell.
\]
The exponent in \eqref{eq:coefficient-extraction} becomes
\(\sum_{\ell=1}^k\zeta_\ell\cdot\eta_\ell\).
We have thus obtained a Walsh test on \(V_{kn'}\), of amplitude
\(\widehat h(P)\), with separate unit-ball families on its two sides.
Monotonicity and amplification invariance now yield
\[
w_{2n'}(T)\ge w_{kn'}(T)
=w_{kn'}(\widetilde T)\ge|\widehat h(P)|
\ge(c_*/2)w_n(T).
\]
Taking \(c_1=\min(1/2,c_*/2)\) proves the lemma.
\end{proof}

\subsection{Uniform decay}

The growth lemma turns a sufficiently late slow decrease into an
infinite sequence of slow decreases for the same operator. For a
finite-rank operator, this would contradict its Hilbert-space bound.
This is the only point at which finite rank is used.

\begin{theorem}[Uniform finite-rank decay]\label{thm:uniform-decay}
Fix \(m,D_g,D_f\) satisfying the assumptions of Lemma~\ref{lem:growth}.
There is an integer \(n_0\), depending only on these parameters,
such that every finite-rank \(T:G\to F\) with
\(c_m(G^*)\le D_g\) and \(c_m(F)\le D_f\) satisfies
\begin{equation}\label{eq:uniform-drop}
w_n(T)\le\tfrac12 w_{n-2m}(T)\qquad(n\ge n_0).
\end{equation}
In particular there are numbers \(\rho_N\to0\), independent of
\(G,F,T\), such that
\begin{equation}\label{eq:uniform-decay}
w_N(T)\le\rho_N\|T\|.
\end{equation}
Both bounds also hold for every \(T:G\to F\) that is an
operator-norm limit of finite-rank maps, under the same moment
hypotheses on \(G^*\) and \(F\).
\end{theorem}
\begin{proof}
Let \(c_1\) be given by Lemma~\ref{lem:growth}. Choose an integer
\(K>\log_2(1/c_1)\), and put
\[
n_0=4m(K+2)+2.
\]
Fix a nonzero finite-rank \(T\), and suppose that
\eqref{eq:slow-decrease} holds at some \(n\ge n_0\).
The growth lemma gives
\[
w_{2n-4m}(T)\ge c_1w_n(T).
\]
By monotonicity, the same lower bound holds throughout
\([n,2n-4m]\).

Set \(a_n=\lceil3n/2\rceil\). The \(K\) consecutive blocks
\[
[a_n+2m(j-1),\,a_n+2mj]\qquad(1\le j\le K)
\]
fit in this interval, since
\(2n-a_n=\lfloor n/2\rfloor\ge2mK+4m+1\).
Their endpoint ratios telescope:
\[
\prod_{j=1}^K
\frac{w_{a_n+2m(j-1)}(T)}{w_{a_n+2mj}(T)}
=\frac{w_{a_n}(T)}{w_{a_n+2mK}(T)}\le\frac1{c_1}.
\]
Since \(2^K>1/c_1\), one right endpoint \(t\) satisfies
\[
w_{t-2m}(T)\le2w_t(T),\qquad
t\ge3n/2,\qquad w_t(T)\ge c_1w_n(T).
\]
We can repeat the argument at \(t\).
Every application of the growth lemma returns an inequality for
the original \(T\); its auxiliary amplifications are discarded.
We therefore obtain indices \(n_j\ge(3/2)^j n\) with
\(w_{n_j}(T)\ge c_1^j w_n(T)>0\).
The rank-dependent estimate \eqref{eq:rank-dependent-decay}, for
this one fixed operator, implies
\[
0<w_n(T)\le D_Tc_1^{-j}2^{-n_j/2}\longrightarrow0,
\]
a contradiction.
Thus \eqref{eq:slow-decrease} is impossible for \(n\ge n_0\), proving
the strict form of \eqref{eq:uniform-drop} for nonzero finite-rank
\(T\). The zero operator satisfies the asserted weak inequality.

Iteration proves \eqref{eq:uniform-decay}; for example, take
\[
\rho_N=
\begin{cases}
1,&N<n_0,\\
2^{-\lfloor(N-n_0)/(2m)\rfloor},&N\ge n_0.
\end{cases}
\]
Finally, if finite-rank maps \(T_j\) converge in operator norm to
\(T\), then \eqref{eq:walsh-continuity} gives
\(w_s(T_j)\to w_s(T)\) for every \(s\), while
\(\|T_j\|\to\|T\|\). Passing to the limit proves both estimates
for \(T\).
\end{proof}

\begin{remark}
The constant \(D_T\) may depend on the operator rank. It is used only
to rule out an infinite sequence of slowly decreasing indices for
one fixed operator; the resulting \(n_0\) and \(\rho_N\) depend only
on the fixed moment parameters. The binary ranks of Fourier indices
in Lemmas~\ref{lem:coefficient} and~\ref{lem:growth} play a different role: they
count the independent Walsh factors extracted from a coefficient.
\end{remark}

\section{Compact factorization and the approximation property}
\label{sec:approximation}

We now transfer Walsh decay to the identity of \(X\).
The approximation property will let us apply Theorem~\ref{thm:uniform-decay}
to compact maps into \(X\). We therefore factor the inclusion of
each finite-dimensional \(E\subset X\) through a compact map
whose domain dual has a controlled moment estimate.
Using \(E\) itself as the domain would not suffice: the quotient
\(E^*\) of \(X^*\) need not inherit the cotype estimate of \(X^*\).
Instead, we enlarge \(E\) to retain the moment lifts available in
\(X\), and use decreasing weights to obtain compactness.

\begin{lemma}[A compact domain with controlled dual moments]
\label{lem:compact-domain}
Let \(m\ge1\) be an integer, let \(\Delta\ge1\), and let \(X\) be
a Banach space with \(c_m(X^*)\le\Delta\).
For every nonzero finite-dimensional subspace \(E\subset X\), there
exist a Banach space \(G\), an injective compact operator \(J:G\to X\),
and a contraction \(\iota:E\to G\), such that
\begin{equation}\label{eq:compact-domain}
J\iota x=x\quad(x\in E),\qquad
\|J\|\le2,\qquad c_m(G^*)\le12\Delta.
\end{equation}
\end{lemma}
\begin{proof}
\textit{Finite-dimensional enlargements.}
We construct nested finite-dimensional subspaces
\[
E=E_0\subset E_1\subset E_2\subset\cdots\subset X
\]
such that each \(z=(z_i)_{i=1}^m\in\ell_2^m(E_l)\) has a lift
\(H_l:\{-1,1\}^m\to E_{l+1}\) satisfying
\begin{equation}\label{eq:finite-level-lift}
\E_\eps\eps_iH_l(\eps)=z_i,\qquad
\|H_l\|_{L_2(E_{l+1})}\le3\Delta\|z\|_{\ell_2^m(E_l)}.
\end{equation}
First take moments in the whole space \(X\). The map
\[
U:L_2(\{-1,1\}^m;X)\to\ell_2^m(X),\qquad
UH=(\E_\eps\eps_iH(\eps))_i,
\]
has the bounded right inverse
\((z_i)\mapsto\sum_i\eps_i z_i\), and
\[
\|U^*(\theta_i)\|_{L_2(X^*)}
=\left(\E_\eps\left\|\sum_i\eps_i\theta_i\right\|^2\right)^{1/2}
\ge\Delta^{-1}\|(\theta_i)\|_{\ell_2^m(X^*)}.
\]
By \eqref{eq:quotient-lift}, the infimal lift norm of a unit tuple
is at most \(\Delta\), so a lift of norm below \(2\Delta\) exists.
Take a finite \(\eta\)-net on the unit sphere of \(\ell_2^m(E_l)\),
where \(0<\eta\le\min(1/2,\Delta/\sqrt m)\), and choose such a lift
for each net point. Let \(E_{l+1}\) be the span of \(E_l\) and all
the values of these finitely many lifts.

For any unit tuple \(z\in\ell_2^m(E_l)\), choose a net point \(v\)
with \(\|z-v\|_{\ell_2^m(X)}\le\eta\). Add
\(\sum_i\eps_i(z_i-v_i)\) to its lift. This corrects the moments
exactly and adds at most
\[
\left\|\sum_i\eps_i(z_i-v_i)\right\|_{L_2(X)}
\le\sqrt m\,\|z-v\|_{\ell_2^m(X)}\le\Delta
\]
to the norm. This proves \eqref{eq:finite-level-lift} for unit
tuples; scaling and the zero lift cover the other cases.

\medskip
\textit{The compact map.}
Put
\[
\mathcal H=\left(\bigoplus_{l=0}^\infty E_l\right)_{\ell_2},
\qquad
j:\mathcal H\to X,\qquad j(z)=\sum_{l\ge0}2^{-l}z_l.
\]
Cauchy--Schwarz proves absolute convergence of the series and gives
\[
\|j\|\le\left(\sum_{l\ge0}4^{-l}\right)^{1/2}<2.
\]
The same estimate on the tails shows that the finite-rank truncations
converge in operator norm. Thus \(j\) is compact.
Let \(G=\mathcal H/\ker j\), with quotient map \(Q\), and define
\(J\) by \(JQ=j\). The kernel is closed, so \(G\) is Banach;
the quotient norm gives \(\|J\|\le\|j\|<2\), and \(J\) is injective.
Every element of \(B_G\) has a representative in \(2B_{\mathcal H}\).
Hence \(J(B_G)\subset j(2B_{\mathcal H})\) is relatively compact,
which proves compactness of \(J\). Finally,
\(\iota(x)=Q(x,0,0,\ldots)\) is a contraction and \(J\iota x=x\).
The finite-rank truncations of \(j\) need not annihilate \(\ker j\),
so they need not induce operators on \(G\). Thus this compactness
argument does not itself approximate \(J\) by finite-rank maps;
that passage will use the approximation property of \(X\).

\medskip
\textit{Dual moments from the weighted shift.}
The nesting of the \(E_l\)'s makes
\[
R:\mathcal H\to\mathcal H,\qquad
R(z_0,z_1,\ldots)=(0,2z_0,2z_1,\ldots)
\]
a bounded map. The factor \(2\) compensates for the next dyadic
weight, so \(jR=j\) and therefore \(QR=Q\).
This identity lets us assemble the finite-level lifts into exact
moment lifts in \(G\).

For \(g_1,\ldots,g_m\in G\), choose representatives
\(z^i=(z_l^i)_l\in\mathcal H\) with
\(Qz^i=g_i\) and \(\|z^i\|\le2\|g_i\|\), taking \(z^i=0\)
when \(g_i=0\).
Use \eqref{eq:finite-level-lift} on each tuple \((z_l^i)_i\),
and set
\[
\widetilde H(\eps)=(0,2H_0(\eps),2H_1(\eps),\ldots).
\]
Summing squared estimates yields
\[
\|\widetilde H\|_{L_2(\mathcal H)}
\le6\Delta\left(\sum_i\|z^i\|^2\right)^{1/2}
\le12\Delta\left(\sum_i\|g_i\|^2\right)^{1/2}.
\]
Every atom of the sign cube has positive mass, so this estimate
also shows that \(\widetilde H(\eps)\in\mathcal H\) for every
\(\eps\). Its \(i\)-th moment is \(Rz^i\). Since \(QR=Q\),
\[
\E_\eps\eps_iQ\widetilde H(\eps)=QRz^i=Qz^i=g_i.
\]
For \(\theta_1,\ldots,\theta_m\in G^*\), it follows that
\begin{align*}
\left|\sum_i\theta_i(g_i)\right|
&=\left|\E_\eps
\left\langle Q\widetilde H(\eps),\sum_i\eps_i\theta_i\right\rangle\right|\\
&\le12\Delta
\left(\sum_i\|g_i\|^2\right)^{1/2}
\left(\E_\eps\left\|\sum_i\eps_i\theta_i\right\|^2\right)^{1/2}.
\end{align*}
Taking the supremum over \(\sum_i\|g_i\|^2\le1\) gives
\(c_m(G^*)\le12\Delta\) by finite \(\ell_2\) duality.
\end{proof}

We now complete the proof of the main theorem.

\begin{proof}[Proof of Theorem~\ref{thm:main}]
First suppose \(K(X)<\infty\).
On each finite cube, the adjoint of the Rademacher projection on
\(L_2(X)\) is the corresponding projection on \(L_2(X^*)\).
Therefore \(K(X^*)=K(X)\).
By the classical characterization recalled in the introduction,
both \(X\) and \(X^*\) have type strictly greater than one.

If a space \(Y\) has type \(p>1\), with constant \(D\), then \(Y^*\)
has cotype \(p'=p/(p-1)\).
Indeed, for finite tuples \(y_i\in Y\), \(y_i^*\in Y^*\),
\begin{align*}
\left|\sum_i y_i^*(y_i)\right|
&=\left|\E_\eps
\left\langle\sum_i\eps_i y_i,\sum_i\eps_i y_i^*\right\rangle\right|\\
&\le D\left(\sum_i\|y_i\|^p\right)^{1/p}
\left(\E_\eps\left\|\sum_i\eps_i y_i^*\right\|^2\right)^{1/2}.
\end{align*}
Taking the supremum over the unit ball of \(\ell_p^m(Y)\) proves
the claim. Apply it to \(Y=X\) and \(Y=X^*\), and restrict the
resulting cotype inequality on \(X^{**}\) to \(X\).
This direction does not use the approximation property.

Conversely, suppose that \(X\) has the approximation property and
\[
C_q(X)\le C_F,\qquad C_r(X^*)\le C_Y
\]
for finite \(q,r\ge2\) and finite constants \(C_F,C_Y\ge1\).
For an integer \(m\), set
\[
D_f=C_Fm^{1/2-1/q},\qquad
\Delta=C_Ym^{1/2-1/r},\qquad D_g=12\Delta.
\]
The parameters in \eqref{eq:lift-parameters} then satisfy
\[
ma=\frac{m^{1/r}}{24LC_Y},\qquad
mb=\frac{m^{1/q}}{2LC_F}.
\]
These positive powers of \(m\) eventually exceed any fixed multiple
of \((\log m)^3\). Choose \(m\) once so that the hypotheses of
Lemma~\ref{lem:coefficient} hold, and fix the decay sequence \(\rho_N\)
given by Theorem~\ref{thm:uniform-decay} for \(m,D_g,D_f\).
All these parameters are independent of the finite-dimensional
subspaces and approximating operators chosen below.

We first show that
\begin{equation}\label{eq:identity-walsh-decay}
w_N(\Id_X)\le2\rho_N\qquad(N\ge0).
\end{equation}
For a nonzero finite-dimensional \(E\subset X\), choose
\(J:G\to X\) and \(\iota:E\to G\) by Lemma~\ref{lem:compact-domain}.
Apply the approximation property to the compact set
\(M=\overline{J(B_G)}\). For every \(j\ge1\), it supplies a bounded
finite-rank map \(S_j:X\to X\) such that
\[
\|S_jJ-J\|
\le\sup_{x\in M}\|S_jx-x\|<1/j.
\]
Thus \(J\) is an operator-norm limit of finite-rank maps from
\(G\) to \(X\). Since \(c_m(G^*)\le D_g\) and \(c_m(X)\le D_f\),
the norm-limit assertion of Theorem~\ref{thm:uniform-decay} gives
\begin{equation}\label{eq:compact-walsh-decay}
w_N(J)\le\rho_N\|J\|\le2\rho_N.
\end{equation}
Only the composites \(S_jJ\) converge in operator norm;
no common bound on \(\|S_j\|\) is required.

Now fix any input \(x=(x_u)_{u\in V_N}\in L_2(V_N;X)\).
If \(x=0\), there is nothing to prove. Otherwise take
\(E=\Span\{x_u:u\in V_N\}\) and choose the corresponding \(J,\iota\).
The family \(z_u=\iota(x_u)\) satisfies
\(\|z\|_{L_2(G)}\le\|x\|_{L_2(X)}\) and \(Jz_u=x_u\).
Consequently,
\[
\|W_N^{\Id_X}x\|_{L_2(X)}
=\|W_N^Jz\|_{L_2(X)}
\le2\rho_N\|x\|_{L_2(X)},
\]
which proves \eqref{eq:identity-walsh-decay}.

If \(K(X)=\infty\), the characterization
\eqref{eq:classical-characterization} supplies \(\lambda\ge1\)
such that, for every \(N\ge1\), there are \(x_u\in B_X\), \(u\in V_N\),
with
\[
\left\|\sum_{u\in V_N}\alpha_u x_u\right\|
\ge\lambda^{-1}\sum_{u\in V_N}|\alpha_u|
\]
for all real coefficients. These vectors form an input of \(L_2\)
norm at most one, while every Walsh output satisfies
\[
\|(W_N^{\Id_X}x)(v)\|
=\left\|2^{-N}\sum_{u\in V_N}\chi(u,v)x_u\right\|
\ge\lambda^{-1}.
\]
Thus \(w_N(\Id_X)\ge\lambda^{-1}\) for every \(N\), contradicting
\eqref{eq:identity-walsh-decay} and \(\rho_N\to0\).
This proves \(K(X)<\infty\).
\end{proof}

\begin{proof}[Proof of Corollary~\ref{cor:uniform-cotype-entropy}]
If no uniform \(\kappa\) existed, choose spaces \(X_j\) with the
stated fixed cotype bounds and AP, but \(K(X_j)\to\infty\), and set
\[
Z=\left(\bigoplus_{j\ge1}X_j\right)_{\ell_2}.
\]
Finite-coordinate truncations converge uniformly on each compact subset
of \(Z\). Approximating its finitely many compact coordinate images by
finite-rank maps therefore proves AP for \(Z\); no uniform bound on
the norms of those maps is needed. Moreover,
\(Z^*=(\bigoplus_{j\ge1}X_j^*)_{\ell_2}\) isometrically.
For any finite tuple \(z_i=(z_{ij})_j\) in \(Z\), Minkowski's inequality
in \(\ell_{q/2}\) and the coordinate cotype bounds give
\[
\left(\sum_i\|z_i\|_Z^q\right)^{2/q}
\le \sum_j\left(\sum_i\|z_{ij}\|^q\right)^{2/q}
\le C^2\E_\eps\left\|\sum_i\eps_i z_i\right\|_Z^2.
\]
Thus \(C_q(Z)\le C\), and the same argument on the dual gives
\(C_r(Z^*)\le D\). Theorem~\ref{thm:main} implies \(K(Z)<\infty\).
Restriction of each finite-cube Rademacher projection to a coordinate
summand gives \(K(X_j)\le K(Z)\), a contradiction.

For the entropy assertion, finite-dimensional gauge spaces have AP,
so the first assertion bounds their \(K\)-constants by \(\kappa\).
The bounded-\(K\) entropy-duality theorem of
Artstein--Milman--Szarek--Tomczak-Jaegermann \cite[Theorem~5]{AMSTJ2004},
in the covering-number formulation of
\cite[Theorem~1.1, equation~(1.6), and Section~2]{Milman2007},
has constants \(a,b\ge1\) depending only on an upper bound for
the \(K\)-convexity constant of either gauge space.
Polarity identifies \(X_{B^\circ}=X_B^*\), and the finite-cube
adjoint identity used above gives \(K(X_B^*)=K(X_B)\).
Thus the same bound is available for the polar bodies in the reverse
inequality. Also, \(X_{tL}\) is linearly isometric to \(X_L\).
Apply the covering theorem to \(K,tL\). Since
\((tL)^\circ=t^{-1}L^\circ\) and
\(N(t^{-1}L^\circ,cK^\circ)=N(L^\circ,ctK^\circ)\) for \(c>0\),
both inequalities in \eqref{eq:cotype-entropy-duality} follow for every
\(t>0\), with constants depending only on \(\kappa\).
\end{proof}

\section{The conjecture on Grothendieck pairs}\label{sec:grothendieck}

We now apply the characterization to the Grothendieck-pair conjecture.
The deduction through complemented Euclidean subspaces is the one in
\cite[Corollary 4.8]{GMR25}; the main theorem allows ordinary AP here.
For an algebraic tensor \(t\in E\otimes F\), recall the injective
and projective norms
\[
\|t\|_\eps
=\sup_{\substack{x^*\in B_{E^*}\\y^*\in B_{F^*}}}
  |(x^*\otimes y^*)(t)|,\qquad
\|t\|_\pi
=\inf_{t=\sum_i x_i\otimes y_i}\sum_i\|x_i\|\|y_i\|.
\]
Their completions are denoted by
\(E\widehat\otimes_\eps F\) and \(E\widehat\otimes_\pi F\).
Both norms satisfy the metric mapping property:
tensoring bounded maps multiplies the norm by at most the product
of their operator norms. This follows directly from the two
definitions.

Let \(\mathcal E=(E_n)_{n\ge1}\) be a sequence of real Banach spaces
with \(\dim E_n=n\).
Following \cite[Definition 1.2]{GMR25}, call \((\mathcal E,F)\)
a \emph{Grothendieck pair} if there is \(C<\infty\) such that
\begin{equation}\label{eq:grothendieck-pair}
\|A\otimes\Id_F:
E_m\widehat\otimes_\eps F
\longrightarrow E_n^*\widehat\otimes_\pi F\|
\le C\|A\|
\end{equation}
for all \(m,n\ge1\) and every \(A:E_m\to E_n^*\).
The existence of the bounded extension in
\eqref{eq:grothendieck-pair} is part of the condition.

\begin{corollary}\label{cor:grothendieck}
If \((\mathcal E,F)\) is a Grothendieck pair and \(F\) is
infinite-dimensional with the approximation property, then
\((\mathcal E,\ell_2)\) is a Grothendieck pair.
In particular, Pisier's conjecture
\cite[Conjecture 1.4]{GMR25} holds, with the approximation property
in place of the bounded approximation property.
\end{corollary}
\begin{proof}
By Gupta--Misra--Ray \cite[Theorem 1.3]{GMR25}, the pair condition
implies finite cotype of \(F\) and \(F^*\).
Thus Theorem~\ref{thm:main} makes \(F\) \(K\)-convex.

We use the complemented Euclidean subspace theorem of
Figiel--Tomczak-Jaegermann \cite{FTJ79}, in the formulation
\cite[Theorem~4]{Pis84ICM}: for an infinite-dimensional
\(K\)-convex space \(F\), there is \(D<\infty\) such that for every
\(d\ge1\) there are bounded maps
\[
U_d:\ell_2^d\to F,\qquad V_d:F\to\ell_2^d,\qquad
V_dU_d=\Id_{\ell_2^d},\qquad \|U_d\|\|V_d\|\le D.
\]
By the metric mapping property and the factorization
\[
A\otimes\Id_{\ell_2^d}
=(\Id_{E_n^*}\otimes V_d)
 (A\otimes\Id_F)
 (\Id_{E_m}\otimes U_d),
\]
the corresponding tensor operator with \(\ell_2^d\) in place of
\(F\) has norm at most \(CD\|A\|\), uniformly in \(d,m,n,A\).

Every algebraic tensor in \(E_m\otimes\ell_2\) has its second
coordinates in a finite-dimensional Hilbert subspace \(H\subset\ell_2\).
The inclusion of \(H\) and the orthogonal projection onto \(H\)
both have norm one. The metric mapping property therefore shows
that both tensor norms on tensors supported in \(H\) agree with
the corresponding norms computed in \(\ell_2\).
The uniform estimate for finite-dimensional Hilbert spaces
consequently holds on \(E_m\otimes\ell_2\), and extends to its
injective completion. This is \eqref{eq:grothendieck-pair} for
\(\ell_2\), with constant \(CD\).
\end{proof}

\begingroup
\small
\raggedright
\begin{thebibliography}{99}

\bibitem{AMSTJ2004}
S. Artstein, V. Milman, S. J. Szarek, and N. Tomczak-Jaegermann,
On convexified packing and entropy duality,
\emph{Geom. Funct. Anal.} \textbf{14} (2004), no.~5, 1134--1141.
\href{https://doi.org/10.1007/s00039-004-0486-3}
{doi:10.1007/s00039-004-0486-3}.
\href{https://arxiv.org/abs/math/0407238v1}{arXiv:math/0407238v1};
numbered locators refer to this version.

\bibitem{Borell79}
C. Borell,
On the integrability of Banach space valued Walsh polynomials,
\emph{S\'eminaire de probabilit\'es}, \textbf{13} (1979), 1--3.
\href{https://www.numdam.org/item/SPS_1979__13__1_0/}{Numdam}.

\bibitem{Enf73}
P. Enflo,
A counterexample to the approximation problem in Banach spaces,
\emph{Acta Math.} \textbf{130} (1973), 309--317.
\href{https://doi.org/10.1007/BF02392270}{doi:10.1007/BF02392270}.

\bibitem{FJ73}
T. Figiel and W. B. Johnson,
The approximation property does not imply the bounded approximation property,
\emph{Proc. Amer. Math. Soc.} \textbf{41} (1973), no.~1, 197--200.
\href{https://doi.org/10.1090/S0002-9939-1973-0341032-5}
{doi:10.1090/S0002-9939-1973-0341032-5}.

\bibitem{FTJ79}
T. Figiel and N. Tomczak-Jaegermann,
Projections onto Hilbertian subspaces of Banach spaces,
\emph{Israel J. Math.} \textbf{33} (1979), no.~2, 155--171.
\href{https://doi.org/10.1007/BF02760556}{doi:10.1007/BF02760556}.

\bibitem{GMR25}
R. Gupta, G. Misra, and S. K. Ray,
\emph{On variants of the Grothendieck inequality and estimates on
tensor product norms},
12 August 2025.
\href{https://arxiv.org/abs/2305.13270v2}{arXiv:2305.13270v2}.

\bibitem{MP76}
B. Maurey and G. Pisier,
S\'eries de variables al\'eatoires vectorielles ind\'ependantes et
propri\'et\'es g\'eom\'etriques des espaces de Banach,
\emph{Studia Math.} \textbf{58} (1976), 45--90.
\href{https://doi.org/10.4064/sm-58-1-45-90}{doi:10.4064/sm-58-1-45-90}.

\bibitem{Milman2007}
E. Milman,
A remark on two duality relations,
\emph{Integral Equations Operator Theory} \textbf{57} (2007),
no.~2, 217--228.
\href{https://doi.org/10.1007/s00020-006-1479-4}
{doi:10.1007/s00020-006-1479-4}.
\href{https://arxiv.org/abs/math/0603461v2}{arXiv:math/0603461v2};
numbered locators refer to this version.

\bibitem{Pis80}
G. Pisier,
Un th\'eor\`eme sur les op\'erateurs lin\'eaires entre espaces de Banach
qui se factorisent par un espace de Hilbert,
\emph{Ann. Sci. \'Ecole Norm. Sup.} (4) \textbf{13} (1980), 23--43.
\href{https://doi.org/10.24033/asens.1376}{doi:10.24033/asens.1376}.

\bibitem{Pis80Kseminar}
G. Pisier,
Semi-groupes holomorphes et \(K\)-convexit\'e,
\emph{S\'eminaire d'analyse fonctionnelle, \'Ecole Polytechnique},
1980--1981, expos\'e II, 1--30.
\href{https://www.numdam.org/item/SAF_1980-1981____A2_0/}{Numdam}.

\bibitem{Pis81}
G. Pisier,
Semi-groupes holomorphes et \(K\)-convexit\'e (suite),
\emph{S\'eminaire d'analyse fonctionnelle, \'Ecole Polytechnique},
1980--1981, expos\'e VII, 1--9.
\href{https://www.numdam.org/item/SAF_1980-1981____A7_0/}{Numdam}.

\bibitem{Pis82K}
G. Pisier,
Holomorphic semi-groups and the geometry of Banach spaces,
\emph{Ann. of Math.} (2) \textbf{115} (1982), no.~2, 375--392.
\href{https://doi.org/10.2307/1971396}{doi:10.2307/1971396}.

\bibitem{Pis83}
G. Pisier,
Counterexamples to a conjecture of Grothendieck,
\emph{Acta Math.} \textbf{151} (1983), 181--208.
\href{https://doi.org/10.1007/BF02393206}{doi:10.1007/BF02393206}.

\bibitem{Pis84ICM}
G. Pisier,
Finite rank projections on Banach spaces and a conjecture of Grothendieck,
in \emph{Proceedings of the International Congress of Mathematicians
(Warsaw, 1983)}, Vol.~2, PWN, Warsaw, 1984, 1027--1039.
\href{https://www.imj-prg.fr/wp-content/uploads/2020/prix/pisier1983.pdf}{Text}.

\bibitem{Pis11IHP}
G. Pisier,
\emph{Martingales in Banach Spaces
(in connection with Type and Cotype)},
IHP course notes, version of 9 February 2011.
\href{https://webusers.imj-prg.fr/~gilles.pisier/ihp-pisier.pdf}{Author's text}.

\bibitem{Pis12}
G. Pisier,
Grothendieck's theorem, past and present,
\emph{Bull. Amer. Math. Soc.} (N.S.) \textbf{49} (2012), no.~2,
237--323.
\href{https://doi.org/10.1090/S0273-0979-2011-01348-9}{doi:10.1090/S0273-0979-2011-01348-9}.

\bibitem{Sza87}
S. J. Szarek,
A Banach space without a basis which has the bounded approximation property,
\emph{Acta Math.} \textbf{159} (1987), 81--98.
\href{https://doi.org/10.1007/BF02392555}{doi:10.1007/BF02392555}.

\bibitem{STJ09}
S. J. Szarek and N. Tomczak-Jaegermann,
On the nontrivial projection problem,
\emph{Adv. Math.} \textbf{221} (2009), no.~2, 331--342.
\href{https://doi.org/10.1016/j.aim.2008.12.004}{doi:10.1016/j.aim.2008.12.004}.

\end{thebibliography}
\endgroup

\end{document}
